\begin{abstract}
Converting transformers to sub-quadratic architectures enables efficient inference but discards few-shot adaptation capability---a tradeoff accepted as fundamental. We show this tradeoff stems from sequential conversion-then-adaptation pipelines, not architectural constraints. We introduce Task-Conditioned Selective State Space (TC-SSM), which integrates task conditioning directly into architecture conversion. Our key insight is that task-relevant structure already exists in pretrained hidden states and can be discovered via unsupervised clustering, enabling task conditioning without labels. Low-rank projections modulate Mamba's state space parameters based on learned task embeddings, adding only 1\% parameters while counterintuitively reducing inference time to 0.85$\times$ vanilla Mamba. On SuperGLUE, TC-SSM preserves transformer-level adaptation: 0.25\% accuracy gap versus baseline and adaptation in just 14 gradient steps---20$\times$ and 7$\times$ better than required thresholds. Our results demonstrate that sub-quadratic efficiency and adaptation capability need not be traded off when conversion and conditioning are co-designed.
\end{abstract}
